% ATTEMPT_STATUS: unresolved
% RESEARCH_GENERATED: 2026-10-01; GPT-6 Astra Ultra; individual mathematical attempt
% TOP_PROBLEM: {"id": 125, "title": "Primes with p-2 Having Odd Omega", "category_id": 1, "category": {"id": 1, "name": "number_theory", "display_name": "Number Theory", "description": "Properties of integers, prime numbers, Diophantine equations.", "slug": "number-theory", "order_index": 1, "created_at": "2026-07-31T15:26:25.670Z"}, "status": "open"}
% FIRST_POSTED: 2026-10-01
% Imported from: unsolvedmath_additional_500.tex; UnsolvedMath ID 125
% !TeX program = lualatex
% Compile twice: lualatex 125.tex
% Self-contained: no external bibliography, images, or \input files.
% Numeric section labels and visible numbers are original UnsolvedMath IDs.
\documentclass[11pt,a4paper]{article}
\usepackage[margin=24mm,headheight=14pt]{geometry}
\usepackage{fontspec}
\setmainfont{Latin Modern Roman}
\setsansfont{Latin Modern Sans}
\setmonofont{Latin Modern Mono}
\usepackage{amsmath,amssymb,mathtools}
\usepackage{unicode-math}
\setmathfont{Latin Modern Math}
\usepackage{microtype}
\usepackage{enumitem,xurl,needspace}
\usepackage[dvipsnames]{xcolor}
\usepackage{hyperref}
\hypersetup{unicode=true,colorlinks=true,linkcolor=MidnightBlue,urlcolor=MidnightBlue,
 pdftitle={UnsolvedMath Problem 125},pdfauthor={Source-annotated compilation},bookmarksnumbered=true}
\usepackage{bookmark,tocloft,titlesec,fancyhdr}
\setlength{\cftsecnumwidth}{4.2em}
\cftsetpnumwidth{2.5em}
\cftsetrmarg{4em}
\titleformat{\section}{\Large\bfseries\raggedright}{\thesection}{0.7em}{}
\pagestyle{fancy}\fancyhf{}
\fancyhead[L]{\small\sffamily UnsolvedMath research attempt}
\fancyhead[R]{\small Original catalogue IDs}
\fancyfoot[C]{\thepage}
\setlength{\parindent}{0pt}
\setlength{\parskip}{5pt plus 1pt}
\setlength{\emergencystretch}{3em}
\setcounter{tocdepth}{1}\setcounter{secnumdepth}{1}
\setlist[itemize]{leftmargin=3em,itemsep=4pt,topsep=4pt,parsep=0pt}
\allowdisplaybreaks
\newcommand{\N}{\mathbb{N}}
\newcommand{\Z}{\mathbb{Z}}
\newcommand{\Q}{\mathbb{Q}}
\newcommand{\R}{\mathbb{R}}
\newcommand{\C}{\mathbb{C}}
\newcommand{\F}{\mathbb{F}}
\newcommand{\A}{\mathbb{A}}
\newcommand{\PP}{\mathbb{P}}
\newcommand{\E}{\mathbb{E}}
\newcommand{\eps}{\varepsilon}
\newcommand{\dd}{\,\mathrm d}
\newcommand{\sref}[2]{\hyperlink{source:#1:#2}{[S#2]}}
\newif\ifshowcatalogue
\showcataloguetrue
\newcommand{\cataloguescope}[1]{\ifshowcatalogue
 \paragraph{Statement recorded in the supplied source.}
 #1\fi}
\begin{document}
% UnsolvedMath ID 125; source catalogue code GREEN-037

\setcounter{section}{124}

\section{Prime Shifts with an Odd Number of Prime Factors}\label{125}

{\small\sffamily UnsolvedMath ID 125 · Number Theory}

\subsection{Definitions and mathematical statement}

\paragraph{Shared notation.}Unless a section says otherwise, $\N=\{1,2,\ldots\}$,
$\Z$ is the ring of integers, $\Q,\R,\C$ are the rational, real, and complex
fields, $[N]=\{1,\ldots,\lfloor N\rfloor\}$, and $\log$ is the natural
logarithm. For real functions of a parameter tending to infinity,
$f\ll g$ and $f=O(g)$ mean $|f|\leq Cg$ eventually for some fixed $C$;
$f\gg g$ reverses this comparison; $f=o(g)$ means $f/g\to0$;
$f\sim g$ means $f/g\to1$; and $f\asymp g$ means both order bounds.
A subscript on $O$ or $\ll$ specifies allowed constant dependence.
The expression $n^{o(1)}$ in an upper bound means growth smaller than
$n^\epsilon$ for each fixed $\epsilon>0$. Local definitions govern any
exception, finite-field convention, or use of the same symbol for another
object. An asymptotic claim is not automatically an effective algorithm.



For $m\geq1$, write $m=\prod_p p^{v_p(m)}$ and define $\Omega(m)=\sum_p v_p(m)$, with $\Omega(1)=0$. Equivalently the Liouville function is $\lambda(m)=(-1)^{\Omega(m)}$. The target is infinitely many primes $p\geq3$ with $\lambda(p-2)=-1$; prime factors are counted with multiplicity. \sref{125}{1} \sref{125}{2}

\paragraph{Statement recorded in the supplied source.}
Do there exist infinitely many primes $p$ for which $p-2$ has an odd number of prime factors (counting multiplicity)? \sref{125}{1}

\subsection{Short English statement}

Are there infinitely many primes whose predecessor by two has an odd total number of prime factors, including repetitions?

\subsection{Research attempt}

\paragraph{Provenance and scope.}
This individual attempt was generated by GPT-6 Astra Ultra on 1 October 2026. The method was to derive explicit partial results and test the first proposed extension adversarially. The original statement and sources below are retained. No claim of mathematical novelty or complete solution is made.

\paragraph{Formal target and source audit.}
Write $\lambda(n)=(-1)^{\Omega(n)}$, with multiplicity counted and $\lambda(1)=1$. The target is infinitely many primes $p\geq3$ with $\lambda(p-2)=-1$. The supplied Carella preprint was inspected: its abstract claims strong cancellation for both Mobius and Liouville functions over shifted primes. Those claims would have consequences for the target, but their proofs have not been independently verified in this attempt and are not used as established inputs. The work below gives exact equivalences and an explicit square-divisor expansion, keeping the distinction between the two functions intact.

\paragraph{The exact counting identity.}
For $X\geq3$, set
\[
 P(X)=\#\{3\leq p\leq X:p\text{ prime}\},\quad
 L(X)=\sum_{3\leq p\leq X}\lambda(p-2),\quad
 Q(X)=\#\{3\leq p\leq X:\lambda(p-2)=-1\}.
\]
Since $1_{\lambda=-1}=(1-\lambda)/2$, we have
\[
 Q(X)=\frac{P(X)-L(X)}2.
\]
Thus the desired infinitude is equivalent to unboundedness of $P(X)-L(X)$. A sufficient, stronger condition is that $L(X)\leq(1-\eta)P(X)$ for some fixed $\eta>0$ along an unbounded sequence of $X$. Full cancellation $L(X)=o(P(X))$ would give asymptotic density one half among these primes. This is stronger than required; the proof must not confuse a convenient sufficient statement with an equivalent necessary condition.

If only finitely many suitable primes exist, then eventually $L(X)=P(X)-2Q_0$ for a fixed integer $Q_0$. This makes the first unresolved analytic task particularly crisp: rule out this eventual near-maximal positive correlation. Cancellation of $\lambda$ over all integers alone does not do so, because the primes sample a thin, arithmetically selected set.

\paragraph{Separating Mobius from Liouville exactly.}
For every positive integer $n$,
\[
 \lambda(n)=\sum_{d^2\mid n}\mu(n/d^2).
\]
Both sides are multiplicative, so it suffices to check $n=p^a$. In that case all terms vanish except the final term with remaining exponent zero or one, giving $1$ if $a$ is even and $-1$ if it is odd. This proves the identity, including $n=1$.

Applying it to the shifted-prime sum gives
\[
 L(X)=\sum_{d\leq\sqrt{X-2}}\ 
 \sum_{\substack{3\leq p\leq X\\p\equiv2\pmod{d^2}}}
 \mu((p-2)/d^2).
\]
The $d=1$ term is the shifted Mobius sum. Every $d\geq2$ term is additional, and neither its sign nor its total is determined by that first term. Thus even a verified theorem about $\sum_p\mu(p-2)$ would require an extra argument before being converted into the desired Liouville statement.

\paragraph{A rigorous tail estimate, and the surviving obstacle.}
Truncate the outer sum at $D\geq1$. Bounding $|\mu|\leq1$ and discarding primality, the absolute contribution of $d>D$ is at most
\[
 \sum_{D<d\leq\sqrt X}\left(\frac{X}{d^2}+1\right)
 \ll\frac XD+\sqrt X.
\]
The first term follows by comparison of $\sum_{d>D}d^{-2}$ with an integral. Taking $D=(\log X)^2$ makes this tail $o(X/\log X)$. Consequently cancellation of the finitely growing collection of inner sums for $d\leq D$, with a total error $o(X/\log X)$, would suffice for full Liouville cancellation. This is an actual reduction: its missing input is a uniform family of prime/Mobius correlations in progressions of square modulus, not an unproved interchange of infinite sums. The number of moduli grows with $X$, so individual fixed-modulus estimates with uncontrolled constants are insufficient.

\paragraph{Verification and status.}
The companion exact script factors $p-2$ for every prime up to one thousand. It finds respectively $2,10,82$ qualifying primes up to $10,100,1000$. The prime $3$ contributes $\lambda(1)=1$ and is correctly excluded; prime $2$ is excluded because its shift is zero. These finite counts demonstrate examples but cannot prove infinitude. The square-divisor identity can be checked prime power by prime power and the tail bound uses only nonnegative majorants. The original source's stronger claims have not been promoted to verified theorems. The established output is the exact counting criterion and a controlled truncation reduction. The unresolved step is to prove a nontrivial upper bound on the shifted-prime Liouville correlation, or the stated uniform family of Mobius estimates strong enough to imply one.

\subsection{Sources}

{\small

\begin{itemize}

\item[\hypertarget{source:125:1}{S1}] ULAM.AI, \emph{UnsolvedMath}, supplied \texttt{problems.json}, record 125 (GREEN-037), statement and background. The embedded literature-review date is 2026-08-17; that is the archive’s review date, not a fresh certification. Dataset landing page: \url{https://huggingface.co/datasets/ulamai/UnsolvedMath}.

\item[\hypertarget{source:125:2}{S2}] N. A. Carella, Result on the Mobius Function over Shifted Primes, arXiv:2206.12956 (2022). \url{https://arxiv.org/abs/2206.12956}. Bibliographic information imported from the supplied record.

\item[\hypertarget{source:125:3}{S3}] Ben Green, 100 Open Problems, current Problem 64, updates through December 2025;. \url{https://people.maths.ox.ac.uk/greenbj/papers/open-problems.pdf}. The author-hosted document was inspected on 27 September 2026; the archive’s local code is not a problem-number locator in this paper.

\item[E1] N. A. Carella, \emph{Result on the Mobius Function over Shifted Primes}; abstract inspected, claimed analytic results not independently verified here. \url{https://arxiv.org/abs/2206.12956}. Inspected 1 October 2026.

\end{itemize}

}

\label{end:125}
\end{document}
